%===================================================================%
%===================================================================%
\chapter{Introduction}\label{ch:introduction}
\markboth{Introduction}{}

%===================================================================%
%===================================================================%
%--- Approx Length ---%
% ~5–6 pages
%---------------------%

%===================================================================%
\section{Motivation and Context}
%===================================================================%

% COMMENT: The textile waste statistics below should be verified and cited.
% Candidate sources include the European Environment Agency (EEA) briefing
% on textiles and the environment (2024), the Ellen MacArthur Foundation
% "A New Textiles Economy" report (2017), and McKinsey/Global Fashion
% Agenda reports. Add the appropriate BibTeX entries and replace the
% placeholder citations below. Search for recent EU-level data to ensure
% the numbers are current and defensible.
The global fashion and textile industry generates substantial volumes of post-consumer waste, with a large fraction ending up in landfill or incineration rather than being recovered for reuse or recycling.
% \cite{TODO_EEA2024TextileBriefing,TODO_EllenMacArthur2017NewTextiles}
Within the European Union, regulatory pressure to increase textile recycling rates is growing, with separate collection of textile waste becoming mandatory in member states.
% \cite{TODO_EU_WasteFrameworkDirective2024}
A key bottleneck in the recycling pipeline is the sorting step: incoming textile waste must be classified by material composition, garment type, and condition (e.g., intact, minor damage, major damage) before it can be routed to appropriate downstream processes such as fiber-to-fiber recycling, downcycling, or second-hand resale. Currently, this sorting is performed predominantly by manual labor, which is slow, ergonomically demanding, and difficult to scale.
% \cite{TODO_TextileSortingManualLabor}

Automating textile sorting with robotic systems is a promising approach to increasing throughput and consistency, but it presents significant technical challenges. Garments are deformable objects that exhibit highly variable shapes, complex self-interactions, and contact dynamics that differ fundamentally from the rigid-object manipulation scenarios for which most industrial robots are designed~\cite{Longhini2025ReviewClothManipulation}. Reliable textile sorting requires not only picking individual items from a mixed waste stream but also manipulating them into configurations that allow visual inspection --- for example, stretching a garment flat to expose damage --- and then routing each item to the correct output bin based on the inspection result. This sequence of pick, inspect, and sort naturally involves multiple manipulation primitives and can benefit from coordination between multiple robot arms operating on the same workpiece.

\ac{RL} offers a data-driven approach to learning such complex manipulation behaviors. Rather than manually designing controllers for each manipulation primitive, \ac{RL} enables policies to be discovered through trial-and-error interaction with the task environment, guided by a reward signal that encodes the desired outcome~\cite{SuttonBarto2018,Kober2013Survey}. When extended to \emph{multi-agent} settings, \ac{RL} can learn coordinated policies for teams of robots that must share a workspace, hand off workpieces, and jointly optimize a system-level objective.
% COMMENT: Add citation for multi-agent RL overview here, e.g.,
% \cite{TODO_Busoniu2008MARL} or a more recent survey such as
% \cite{TODO_Zhang2021MARL}. The MARL background section (§2.3) will
% contain the detailed treatment; align the citation once that section
% is written.

Modern \ac{GPU}-accelerated physics simulators have become essential enablers for this approach. By running thousands of environment replicas in parallel, they provide the sample throughput required for stable policy learning while avoiding the cost and risk of training on physical hardware~\cite{Makoviychuk2021IsaacGym}. NVIDIA Isaac~Sim, with its PhysX~5 physics backend and the associated IsaacLab framework, provides an integrated toolchain that combines high-fidelity rigid-body and deformable-object simulation with a modular \ac{MDP}-based task interface suitable for large-scale \ac{RL} training~\cite{nvidia_isaac_sim_510_what_is,Mittal2025IsaacLab}. Recent developments in the Newton physics engine further extend this capability to realistic cloth simulation using Vertex Block Descent (\ac{VBD})~\cite{Chen2024VBD,nvidia_newton_announcement_2025}, making simulation-based development of cloth manipulation policies increasingly practical.

This thesis is the second part of a two-part project conducted at the Institute for Factory Automation and Production Systems (FAPS) at Friedrich-Alexander-Universität Erlangen-Nürnberg. The preceding project thesis (Meyer, 2026) established a validated simulation of the 5-\ac{DoF} tendon-driven robot and demonstrated \ac{RL}-based control for reach and rigid-object place tasks in Isaac~Sim.
% COMMENT: Replace "(Meyer, 2026)" with a proper citation once the
% project thesis is finalized and added to the bibliography, e.g.,
% \cite{Meyer2026PA}.
The present work extends that foundation to the full textile sorting pipeline: integrating deformable cloth simulation, adding a second collaborative robot arm, implementing camera-based damage classification, training coordinated multi-agent policies, and preparing the system for sim-to-real transfer.

%===================================================================%
\section{Problem Statement and Objectives}
%===================================================================%
%---- Problem Statement ----%

Condition-based textile sorting from mixed waste streams poses a multi-faceted robotic manipulation problem that goes substantially beyond standard pick-and-place. The task requires a system that can (i) selectively retrieve target garments (e.g., T-shirts) from a moving conveyor while ignoring non-target objects, (ii) present each garment in a configuration suitable for visual inspection (e.g., stretched flat), (iii) classify the garment's condition through camera-based analysis, and (iv) route the item to the appropriate output bin. Each of these sub-tasks involves contact-rich interaction with deformable objects whose behavior is difficult to model analytically~\cite{Longhini2025ReviewClothManipulation}, and the full pipeline requires coordination between multiple robotic agents operating in a shared workspace.

The robotic system considered in this thesis consists of two manipulators. The first is the tendon-driven 5-\ac{DoF} pick robot validated in the project thesis: a 2-\ac{DoF} prismatic positioning stage, a 3-\ac{DoF} tendon-driven arm, and a Robotiq 2F-140 gripper, based on the design by Klein~\cite{Klein2023}. The second is a collaborative sorting arm that receives the garment from the pick robot, stretches it for camera-based inspection, and deposits it into a condition-specific bin. The two robots must coordinate their actions --- timing the handoff, avoiding collisions, and jointly optimizing the sorting throughput and accuracy --- which naturally frames the problem as a multi-agent \ac{RL} task.

A critical technical challenge is the simulation of cloth dynamics with sufficient fidelity for policy learning. Deformable cloth introduces a high-dimensional state space (many particles/vertices), complex self-collisions and folding behavior, and sensitivity to material parameters such as bending stiffness, friction, and damping~\cite{Longhini2025ReviewClothManipulation,muller2007pbd}. Policies trained on inaccurate cloth models risk exploiting simulation artifacts that do not transfer to real textiles, making the \emph{reality gap} a central concern~\cite{Salvato2021RealityGapSurvey}. Domain randomization over cloth and environment parameters is therefore essential for developing robust policies~\cite{Tobin2017DomainRand,Peng2018DynamicsRandomization}.

%---- Objectives ----%
The thesis pursues the following objectives:

\begin{enumerate}
  \item \textbf{Conveyor-based sorting task --- single agent (PG-1):}
  Extend the validated pick-and-place pipeline from the project thesis to a full sorting task: selective retrieval of target items (T-shirts) from a mixed waste stream on a moving conveyor while ignoring non-target objects. Train \ac{PPO} policies with selectivity encoded in the reward function and evaluate sorting accuracy, false-pick rate, and missed-pick rate.

  \item \textbf{Cloth simulation integration (PG-2):}
  Integrate deformable cloth simulation into the Isaac~Sim environment using PhysX particle cloth (\ac{PBD}) or Newton thin deformables (\ac{VBD})~\cite{muller2007pbd,Chen2024VBD}. Model T-shirt objects with realistic material properties and validate cloth behavior through qualitative and quantitative benchmarks. Apply domain randomization over cloth parameters to support robust policy learning.

  \item \textbf{Second robot arm integration (PG-3):}
  Add a second collaborative robot arm to the simulation scene for cloth stretching and presentation. Define the arm's kinematic configuration, placement, and workspace to complement the tendon-driven picker. Implement the stretching primitive: the second arm grasps a garment edge and extends it to a flat, camera-visible configuration.

  \item \textbf{Camera-based damage classification (PG-4):}
  Integrate a simulated camera system for garment inspection. Develop or integrate a classification model that assesses textile damage degree (e.g., intact, minor damage, major damage, unusable) from the stretched cloth image. Train or fine-tune the classifier on synthetic data with domain randomization.

  \item \textbf{Multi-agent coordination (PG-5):}
  Formulate the full pipeline as a multi-agent \ac{MDP}: the tendon-driven robot picks and hands off, the sorting arm stretches, inspects, and sorts. Design the multi-agent observation and action spaces, inter-agent communication or shared reward structure, and coordination protocol. Train coordinated policies and evaluate end-to-end sorting performance.
  % COMMENT: Decide whether to use Dec-POMDP or centralized-training
  % decentralized-execution (CTDE) framing here. The MARL background
  % section (§2.3) should provide the formal framework. Align terminology
  % once that section is written.

  \item \textbf{Sim-to-real transfer preparation (PG-6):}
  Assess transfer robustness through systematic domain randomization (conveyor speed, cloth properties, sensor noise, lighting)~\cite{Tobin2017DomainRand,Peng2018DynamicsRandomization}. Define the \ac{ROS}~2 integration architecture for physical deployment: joint state publishers, policy inference nodes, camera pipelines, and safety constraints. Identify and document remaining gaps for real-world operation.

  \item \textbf{Evaluation and benchmarking (PG-7):}
  Define and apply comprehensive evaluation metrics: sorting accuracy by damage class, cycle time, false/missed classification rates, cloth coverage during inspection, and multi-agent coordination efficiency. Compare against single-agent baselines and ablation studies (e.g., without stretching, without classification).
\end{enumerate}

%===================================================================%
\section{Thesis Outline}
%===================================================================%

This thesis is structured as follows.
Chapter~\ref{ch:background} provides the background required for this work, beginning with a summary of the foundation established in the project thesis (tendon-driven robot model, simulation environment, and \ac{RL} pipeline) and then covering the additional topics central to the master thesis: cloth and deformable object simulation, multi-agent \ac{RL}, vision-based classification, prior work on robotic cloth manipulation and sorting, and sim-to-real transfer strategies.

Chapter~\ref{ch:research_gap} identifies the specific research gap: the absence of an integrated, simulation-based multi-agent system that combines tendon-driven manipulation, deformable cloth handling, camera-based condition classification, and coordinated sorting within a single \ac{RL}-trainable environment.

Chapter~\ref{ch:methodology} describes the methodological approach, including the cloth simulation setup, the second robot arm integration, the damage classification pipeline, the multi-agent \ac{MDP} formulation, the \ac{RL} training procedures, and the domain randomization strategy for sim-to-real robustness.

Chapter~\ref{ch:results} presents the experimental results across all objectives: single-agent sorting performance, cloth simulation fidelity, multi-agent coordination metrics, classification accuracy, and the end-to-end sorting evaluation.

Chapter~\ref{ch:discussion} discusses the implications and limitations of the approach, including cloth simulation fidelity and its effect on policy transfer, the effectiveness of multi-agent coordination versus sequential single-agent execution, and the remaining challenges for real-world deployment.

Chapter~\ref{ch:conclusion} concludes the thesis, summarizes the contributions, and identifies directions for future work, including physical deployment via the \ac{ROS}~2 integration pathway.
